\documentclass{article}
\usepackage{amsmath,amssymb,amsthm,amsfonts}
\usepackage{geometry}
\usepackage{hyperref}
\usepackage{enumitem}
\geometry{a4paper, margin=1in}

\newtheorem{thm}{Theorem}
\newtheorem{lemma}{Lemma}
\newtheorem{cor}{Corollary}
\newtheorem{exmp}{Example}
\newtheorem{ex}{Exercise}
\newtheorem{defn}{Definition}

\newenvironment{sol}
  {\renewcommand{\qedsymbol}{}\begin{proof}[Solution]}
  {\end{proof}}

\begin{document}

\section{Summatory Function}

\begin{defn}[Summatory Function]
Let $f$ be an arithmetic function. The summatory function of $f$, denoted $S_f(x)$, is the real function defined as:
$$
S_f(x) = 
\begin{cases} 
0 & \text{if } x < 1 \\ 
\sum\limits_{1 \le n \le x} f(n) & \text{otherwise} 
\end{cases}
$$
\end{defn}

\begin{defn}[Bernoulli Polynomials \& Numbers]
The Bernoulli polynomials, $B_n(x)$, are polynomials defined by the generating function as:
\begin{equation}\label{bern:gen:func}
\frac{t e^{tx}}{e^t - 1} = \sum_{n=0}^{\infty} B_n(x) \frac{t^n}{n!}, \quad |t| < 2\pi
\end{equation}
The Bernoulli numbers, $B_n$, are defined as:
$$
\forall n \ge 0, \quad B_n = B_n(0)
$$
\end{defn}

\begin{thm}[Recursive Relation]
	The Bernoulli polynomails satisfies the recursive relation
	$$
	\sum\limits_{k = 1}^{n} \binom{n + 1}{k}B_k(x) = \left( n + 1 \right)x^n
	$$
\end{thm}

\begin{proof}
	By multipling both sides of \eqref{bern:gen:func} by $(\frac{e^t - 1}{t})$ and writing everything as sums and equality of the coefficients the result follows.
\end{proof}

From this, we can compute the first few polynomials:
\begin{align*}
B_0(x) &= 1 \\
\binom{2}{0}B_0(x) + \binom{2}{1}B_1(x) &= 2x \implies B_1(x) = x - \frac{1}{2} \\
\binom{3}{0}B_0(x) + \binom{3}{1}B_1(x) + \binom{3}{2}B_2(x) &= 3x^2 \implies B_2(x) = x^2 - x + \frac{1}{6}
\end{align*}

\medskip

\begin{thm} The Bernoulli polynomails satisfies the following formulas:
\begin{enumerate}
    \item $B_n'(x) = n B_{n-1}(x)$
    \item $B_n(x+1) - B_n(x) = n x^{n-1}$
    \item $B_n(1-x) = (-1)^n B_n(x)$
\end{enumerate}
\end{thm}

\begin{proof}
Follows from the generating function.
\end{proof}

\smallskip

\begin{defn}[Bernoulli Functions]
The Bernoulli functions are the periodic functions with period 1 defined as:
$$
P_n(x) = B_n(\{x\}) = B_n(x - \lfloor x \rfloor)
$$
\end{defn}

\begin{thm}[Euler-Maclaurin Summation Formula]
Let $f(x)$ be a $p$-times continuously differentiable function over $[m, n]$, where $m, n \in \mathbb{N}$ are positive integers with $m < n$. Then for any such $m, n$:
\begin{align*}
\sum\limits_{j=m}^{n} f(j) &= \int\limits_{m}^{n} f(x) dx + \frac{f(m) + f(n)}{2} \\
&+ \sum_{k=1}^{\lfloor p/2 \rfloor} \frac{B_{2k}}{(2k)!} \left( f^{(2k-1)}(n) - f^{(2k-1)}(m) \right) \\
&+ (-1)^{p-1} \int_{m}^{n} \frac{B_p(\{x\})}{p!} f^{(p)}(x) dx
\end{align*}
This formula is called the Euler-Maclaurin summation formula up to the order $p$.
\end{thm}

\begin{proof}
	By induction on the order p the proof follows.
\end{proof}

\smallskip

\begin{exmp} Find the following sum:
	$$
	\sum_{n=1}^{N} \frac{1}{n}
	$$
\end{exmp}

\begin{sol}
We apply the Euler-Maclaurin formula for $f(x) = \frac{1}{x}$ up to order 1:
\begin{align*}
\sum_{n=1}^{N} \frac{1}{n} &= \int_{1}^{N} \frac{1}{x} dx + \frac{1/N + 1}{2} + \int_{1}^{N} \frac{B_1(\{x\})}{1!} \left(\frac{-1}{x^2}\right) dx \\
			   &= \log N + \frac{N+1}{2N} - I_1
\end{align*}
Now we need to calculate the value of the integral $I_1$.
\begin{align*}
	I_1 &= \int_{1}^{N} \frac{\{x\} - 1/2}{x^2} dx \\
	  &= \int_{1}^{N} \frac{\{x\}}{x^2} dx - \frac{1}{2} \int_{1}^{N} \frac{1}{x^2} dx \\
	  &= I_2 + \frac{1 - N}{2N}
\end{align*}
For $I_2$ for $x \in [1, \infty), \, \text{we have that } \dfrac{\{x\}}{x^2} \le \dfrac{1}{x^2}$, we have:
\begin{align*}
	I_2 &= \int_{1}^{N} \frac{\{x\}}{x^2} dx \\
	    &\le \int_{1}^{N} \frac{1}{x^2} dx \\
	    &= \int_{1}^{\infty} \frac{1}{x^2} dx - \int_{N}^{\infty} \frac{1}{x^2} dx \\
	    &\le 1 + \frac{1}{N}
\end{align*}
So $I_2$ is convergant and $I_2 = O(\frac{1}{N})$.

Now collecting everything we get:
$$
\sum_{n=1}^{N} \frac{1}{n} = \log N + \gamma + O\left(\frac{1}{N}\right)
$$
Where $\gamma = \lim_{N \to \infty} \left( \sum_{n=1}^{N} \frac{1}{n} - \log N \right)$.

This constant is called the Euler constant.
\end{sol}

\medskip

\section{Exercises}

\begin{ex}
For $x > 1$, evaluate:
$$
    \sum_{\substack{n \le x \\ n \text{ is odd}}} \frac{1}{n}
$$
\end{ex}

\begin{ex}
    Show that if $k$ is a positive integer, then:
    $$
    \sum_{n \le x} \frac{(\log n)^k}{n} = \frac{(\log x)^{k+1}}{k+1} + C_k + O_k\left(\frac{(\log x)^k}{x}\right)
    $$
    for $x \ge 1$, where $C_k$ is a constant.

\end{ex}

\begin{thm}[Euler's Summation Formula] For $0 < y < x$:
    \begin{align*}
    \sum_{y < n \le x} f(n) &= \int_{y}^{x} f(t) dt + \int_{y}^{x} (t - \lfloor t \rfloor) f'(t) dt \\
    &\quad + f(x)(\lfloor x \rfloor - x) - f(y)(\lfloor y \rfloor - y)
    \end{align*}
    This formula is called the Euler's Summation Formula.
\end{thm}

\begin{proof}
	Apply Euler Maclurain summation formula with order 1.
\end{proof}

\begin{ex}
Show that:
$$
    \sum_{2 \le n \le x} \frac{1}{n \log n} = \log \log x + B + O\left(\frac{1}{x \log x}\right)
$$
\end{ex}

\begin{ex}[Analytical Continuation of Zeta]
	We want to extend the Zeta function by using Euler Macluarin summations.
	\begin{enumerate}[label=(\alph*)]
        \item By applying the Euler-Maclaurin Summation formula to $f(x) = \dfrac{1}{x^s}$, $\text{Re}(s) > 1$, show that $\zeta(s)$ admits an analytic continuation to the half plane $\text{Re}(s) > 1-2k$, except for a simple pole at $s=1$.
        \item Deduce that $\zeta(s)$ can be analytically continued to the entire complex plane $\mathbb{C} \setminus \{1\}$ and that $\zeta(-m) = \dfrac{-B_{m+1}}{m+1}$ holds for any non-negative integer $m$.
        \item Show that $B_{2k+1} = 0$ for all integers $k \ge 1$.
        \item Deduce that $\zeta(-2k) = 0$ for all positive integers $k$.
	\end{enumerate}
\end{ex}

\begin{ex}
    Let $k$ be an arbitrary positive integer. For any $k \ge 2$, find an asymptotic formula for:
    $$
    \sum_{\substack{n \le x \\ \gcd(n, k) = 1}} \frac{1}{n}
    $$
    \textit{Hint:} use the characteristic fuctions.
\end{ex}

\end{document}
